% Generated draft: requires mathematical and source review.
\documentclass[11pt]{article}
\usepackage[margin=25mm]{geometry}
\usepackage{fontspec}
\setmainfont{Latin Modern Roman}
\usepackage{amsmath,amssymb,mathtools,unicode-math}
\setmathfont{Latin Modern Math}
\usepackage{enumitem,needspace,xurl,hyperref}
\setlength{\parindent}{0pt}
\setlength{\parskip}{5pt}
\setlength{\emergencystretch}{3em}
\newcommand{\N}{\mathbb{N}}
\newcommand{\Z}{\mathbb{Z}}
\newcommand{\Q}{\mathbb{Q}}
\newcommand{\R}{\mathbb{R}}
\newcommand{\C}{\mathbb{C}}
\newcommand{\F}{\mathbb{F}}
\newcommand{\A}{\mathbb{A}}
\newcommand{\PP}{\mathbb{P}}
\newcommand{\E}{\mathbb{E}}
\newcommand{\eps}{\varepsilon}
\newcommand{\dd}{\,\mathrm{d}}
\begin{document}
\section*{UnsolvedMath ID: 800413}
\textbf{Source problem code:} catalogue entry
\subsection{Definitions and mathematical statement}
Let \(\mathbb C_v\) denote the ambient valued field appearing in the supplied source, with its \(v\)-topology. Let \(K\subseteq\mathbb C_v\) be a subfield that is locally compact in the induced topology. Let \(\overline K\subseteq\mathbb C_v\) be an algebraic closure of \(K\), embedded over \(K\), and assume that
\[
\operatorname{cl}_{\mathbb C_v}(\overline K)=\mathbb C_v.
\]
Here \(\operatorname{cl}_{\mathbb C_v}\) denotes topological closure. Let \(K(z)\) be the field of rational functions in one variable over \(K\). A rational function \(f\in K(z)\) induces a self-map of \(\mathbb P^1(\mathbb C_v)\); if \(f=P/Q\) with coprime polynomials \(P,Q\in K[z]\), its degree is \(\max\{\deg P,\deg Q\}\).

The conclusion uses the source's notion of a wandering domain. Under the usual interpretation in terms of Fatou components, fix the intended analytic dynamical space and its Fatou set \(\mathcal F_v(f)\). For a Fatou component \(U\), let \(U_n\) be the Fatou component containing \(f^n(U)\), where \(f^0\) is the identity and \(f^n\) is the \(n\)-fold iterate for \(n\geq1\). A wandering domain is a component whose forward component itinerary is not eventually periodic.

\textbf{Conjecture 36 (as supplied).} For every \(K\subseteq\mathbb C_v\) satisfying the local compactness and density hypotheses above, and every \(f\in K(z)\) with \(\deg(f)\geq2\), the map \(f\) has no wandering domains. In the Fatou-component interpretation, this says that for every Fatou component \(U\), there exist integers \(m\geq0\) and \(p\geq1\) such that
\[
U_{n+p}=U_n\qquad\text{for every integer }n\geq m.
\]
\textbf{Unresolved qualification.} The supplied excerpt does not define \(\mathbb C_v\), specify its valuation or characteristic, or identify the analytic space and notion of domain used in the conclusion. In particular, it does not specify a classical or Berkovich interpretation, or whether another non-Archimedean domain convention is intended. The component formulation above is conditional on the intended convention; the excerpt alone does not determine a fully specified analytic formulation.
\subsection{Short English statement}
The supplied conjecture asserts that a rational map of degree at least two over a locally compact subfield \(K\subseteq\mathbb C_v\), whose embedded algebraic closure is dense in \(\mathbb C_v\), has no wandering domains. The ambient field and the precise analytic meaning of a domain require clarification from the source context.
\subsection{Sources}
\begin{itemize}
\item Local UnsolvedMath registry: permanent ID \(800413\), exact source problem code \texttt{catalogue entry}, title ``Open Problems in Arithmetic Dynamics --- Conjecture 36.'' The supplied registry labels this record \texttt{open}; this is supplied curation, not an independently verified mathematical status.
\item Benedetto, DeMarco, Ingram, Jones et al., \emph{Current Trends and Open Problems in Arithmetic Dynamics} (2018), Conjecture 36. Supplied source URL: \url{https://arxiv.org/abs/1806.04980}. Attribution and proposed year are retained from the registry.
\item The source statement includes the reference marker \texttt{[[MR1781331]]}. No corresponding bibliographic entry is included in the supplied record.
\item Supplied provenance records access on 2026-10-16 and extraction by \texttt{arxiv\_latex}. These are source-record metadata; no independent source visit or verification is asserted here.
\end{itemize}
\end{document}
